% !TEX program = xelatex
% 专题讲解入口：先定义概念，再讲机制，最后连接到任务。
\documentclass[aspectratio=169,10pt,fontset=none]{ctexbeamer}
\usetheme{BITUniversal}
% 所有入口共用。复制项目后，先改这里的标题、姓名、机构和日期。
\title[组会汇报]{组会汇报}
\subtitle{研究主题 / 论文中文标题}
\setenglishtitle{Research Topic / Original Paper Title}
\setpresenter{汇报人姓名}
\setaffiliation{北京理工大学}
\setmeeting{课题组组会}
\date{\number\year.\ifnum\month<10 0\fi\number\month.\ifnum\day<10 0\fi\number\day}
% \date{2026.11.09}  % 如需固定日期，取消本行注释。
\setpaper{Original Paper Title}{Author A, Author B, et al.}{Conference / Journal, Year}

% 默认：DT 参考中的轻量白底页眉；改成 band 使用 PWBC 绿色渐变页眉。
\setheaderstyle{light}
% \bitsectionpagesfalse  % 短汇报可关闭自动章节过渡页。

\begin{document}
\BITtitlepage
\BITcontents
\bitsection{问题与基础概念}{应用背景、必要术语与直观理解}
\begin{frame}{应用背景与核心问题}
  \bitimageplaceholder[.44\textheight]{插入观众熟悉的应用场景}
  \bitcallout{问题：\bitplaceholder{这个场景为什么需要本次介绍的方法。}}
\end{frame}
\begin{frame}{概念定义与简单例子}
  \bitcolumns{\begin{block}{定义与符号}
    \bitplaceholder{一句话定义，并给出必要公式。}\end{block}}%
    {\begin{block}{具体例子}
    \bitplaceholder{用一个简单输入说明输出是什么。}\end{block}}
\end{frame}
\bitsection{原理与计算过程}{结构、关键操作与计算步骤}
\begin{frame}{结构与信息流}
  \bitimageplaceholder[.54\textheight]{插入模型或算法结构图}
  \bitsource{来源：原始论文、教材或官方说明。}
\end{frame}
\begin{frame}{关键计算与直观解释}
  \begin{block}{数学表达}
    \[ y_t=f_{\theta}(x_{\leq t},c) \]
  \end{block}
  \bitcallout{直观解释：\bitplaceholder{从输入、计算操作与输出三个方面说明。}}
\end{frame}
\bitsection{应用与讨论}{连接控制任务、比较方法与适用条件}
\begin{frame}{任务建模与训练推理}
  \bitcolumns{\begin{block}{任务如何表示}
    \bitplaceholder{状态、条件、输出的对应关系。}\end{block}}%
    {\begin{block}{模型如何使用}
    \bitplaceholder{训练目标与执行过程。}\end{block}}
\end{frame}
\begin{frame}{总结与适用条件}
  \begin{itemize}
    \item \bitplaceholder{方法的主要能力。}
    \item \bitplaceholder{成立的前提和使用成本。}
    \item \bitplaceholder{与本课题相关的应用方向。}
  \end{itemize}
\end{frame}
\BITthankyou
\end{document}
